\documentclass[a4paper,12pt]{article}
\usepackage[slovene]{babel}
\usepackage[utf8]{inputenc}
\usepackage[T1]{fontenc}
\usepackage{lmodern}
\usepackage{amsmath,amsfonts,amssymb,amsthm}
\usepackage{fancyhdr}
\usepackage{url}
\usepackage{graphicx}
\usepackage{xcolor}

\theoremstyle{definition}
\newtheorem{definicija}{Definicija}

\title{Metric dimension of $C_n(1, 2, \dots, n - d)$: exact values}
\author{Živa Hegler, Neža Zupančič, Sara Polajnar}

\begin{document}
\maketitle

\section{Uvod in definicije}
Pri izvedbi našega projekta si bomo pomagale z definicijami, predstavljenimi
v nadaljevanju. V vseh bomo predpostavile, da imamo usmerjen graf
$G = (V, E)$, kjer je $V$ množica vseh vozlišč in $E$ množica vseh usmerjenih povezav. 
\begin{definicija}
    Naj bodo $1 \leq a_1 < a_2 \ldots < a_k \leq n-1$ generatorji.
    Usmerjeni cirkulantni graf $C_n(a_1, \ldots, a_{k})$ ima vozlišča
    $v_0, v_1, \ldots, v_{n-1}$ in usmerjene povezave $v_iv_{i+a_j}$ za  
    $0 \leq i \leq n-1$ in $1 \leq j \leq k$, pri čemer indekse računamo
    po modulu n. V našem projektu bomo obravnavale usmerjene cirkulantne grafe 
    $C_n(1,2,\ldots,n-d)$, kjer je $n$ število vozlišč in $n-d$ število generatorjev.
\end{definicija}

\begin{definicija}
    Naj bo $W = \{w_1, \ldots, w_z\} \subseteq V$ množica vozlišč.
    Za vsako vozlišče $v \in V$ definiramo vektor razdalj glede na
    množico $W$ kot:
    $$
    r(W|v) = (d(w_1, v), \ldots, d(w_z, v)),
    $$
    kjer $d(w_i, v)$ označuje najkrajšo dolžino usmerjen poti med
    vozliščem $w_i$ in vozliščem $v$. Množiva $W$ je \textbf{razločevalna množica},
    če za vsaki različni vozlišči $u, v \in V$ velja:
    $$
    r(W|u) \neq r(W|v)
    $$
\end{definicija}

\begin{definicija}
    \textbf{Metrična dimenzija} grafa $G$ je moč najmanjše razločevalne
    množice. Označimo jo z $dim(G)$
\end{definicija}

\section{Načrt reševanja problema}

Projekt bomo izvedle v štirih fazah:

\subsection{Naloga 1: Implementacija in verifikacija ILP modela}
Iskanje minimalne razrešujoče množice bomo formulirale kot problem celoštevilskega linearnega programiranja (ILP) 
in pri tem prilagodile obstoječo kodo Maše Popovič. Ker je $n > 2d$, sta možni razdalji med različnimi vozlišči le 
$1$ ali $2$. Z binarnimi spremenljivkami $x_w \in \{0, 1\}$ za vsako vozlišče minimiziramo $\sum_{w \in V} x_w$ pri 
pogojih pokritja
\[
\sum_{w \in D_{u,v}} x_w \geq 1 \quad \text{za vse } u \neq v,
\]
kjer je $D_{u,v} = \{w \in V \mid d(w, u) \neq d(w, v)\}$. Pravilnost kode bomo potrdile z znanima formulama za 
$d = 4$ in $d = 5$ ter s preverjanjem unikatnosti dobljenih vektorjev razdalj.

\subsection{Naloga 2: Računanje za večje grafe in odprava simetrij}
Izvedle bomo sistematične izračune za $4 \leq d \leq 15$ in $n > 2d$, pri čemer bomo za $5 \leq d \leq 10$ vrednosti 
$n$ potisnile vsaj do $100$. Ker reševalniki pri večjih $n$ postanejo počasni, bomo rotacijsko simetrijo cirkulantnega 
grafa odpravile s fiksiranjem začetnega vozlišča ($x_0 = 1$), po potrebi pa preizkusile zmogljivejši zunanji ILP 
reševalnik. Vse dimenzije in čase izvajanja bomo sproti beležile v datoteke CSV.

\subsection{Naloga 3: Testiranje Domneve 1}
Na pridobljenih rezultatih bomo preverile veljavnost Domneve 1, ki postavlja spodnjo mejo:
\[
\dim(C_n(1, \dots, n-d)) \geq B(n, d) = 
\begin{cases}
\left\lceil \frac{n}{d+1} \right\rceil \left\lceil \frac{d+1}{3} \right\rceil, & d \not\equiv 0 \pmod 3, \\[2mm]
\left\lceil \frac{n}{2d-1} \cdot \frac{2d}{3} \right\rceil, & d \equiv 0 \pmod 3.
\end{cases}
\]
Preučile bomo robne primere in pogoje enakosti (zlasti $(d+1) \mid n$ oz. $(2d-1) \mid n$) ter analizirale morebitna 
odstopanja od spodnje meje.

\subsection{Naloga 4: Iskanje eksaktnih formul za $d \geq 6$}
Za fiksne parametre $d \geq 6$ bomo poskusile izpeljati zaprte formule glede na ostanke $n \bmod (d+1)$ oziroma 
$n \bmod (2d-1)$. Določile bomo prag $N_d$, pri katerem se periodični vzorec ustali, in popisale vsa začetna odstopanja 
za majhne $n$.

\end{document}